% !TeX encoding = UTF-8
\documentclass[variant=docente,cover=institutional,mode=screen]{ua-engineering-v6-article}
% !TeX encoding = UTF-8
\UASetUniversity{Universidad Austral}
\UASetFaculty{Facultad de Ingeniería}
\UASetDepartment{Departamento de Computación}
\UASetCampus{Pilar}
\UASetCareer{Ingeniería Informática}
\UASetCommission{Comisión A}
\UASetProfessor{Equipo docente}
\UASetSemester{Segundo cuatrimestre 2026}
\UASetCourse{Sistemas Distribuidos}
\UASetDate{2026}


\UASetClassNumber{10}
\UASetClassTitle{Observabilidad distribuida}
\UASetDocKind{Guía docente}

\title{Clase 10 --- Guía docente: observabilidad, tracing, métricas y debugging}
\author{\UAProfessor}
\date{\UADate}

\begin{document}
\maketitle

\section{Objetivo pedagógico profundo}
La clase debe cambiar la forma en que el estudiante interpreta una falla. El objetivo no es que memorice tres ``pilares'', sino que aprenda a distinguir síntoma, evidencia e hipótesis causal. Al finalizar debería poder explicar por qué una gráfica detecta, una traza localiza y un log confirma detalles, y también por qué cualquiera de esas señales puede faltar, sesgarse o fallar.

\begin{uakeybox}
Idea de cierre: observabilidad no significa acumular telemetría; significa conservar evidencia suficiente para formular y refutar hipótesis sobre el comportamiento del sistema.
\end{uakeybox}

\section{Resultados de aprendizaje esperados}
El estudiante debería poder:
\begin{itemize}
\item diseñar logs estructurados con contexto y sin filtrar secretos;
\item seleccionar Golden Signals, RED y USE según la pregunta;
\item interpretar percentiles y detectar cardinalidad peligrosa;
\item construir una traza sincrónica y otra asincrónica;
\item proponer una investigación causal basada en hipótesis;
\item reconocer que el pipeline de telemetría también necesita resiliencia y SLO propios.
\end{itemize}

\section{Preparación previa del docente}
Preparar antes del encuentro:
\begin{itemize}
\item tres capturas o bocetos del mismo incidente: dashboard, traza y logs;
\item un JSON de log bueno y otro deficiente;
\item una traza de checkout con un span lento y un retry;
\item una tabla pequeña con media, p50, p95 y p99;
\item un caso donde el primer span rojo no sea la causa inicial;
\item una mini-historia inspirada en incidentes reales de pérdida de logs o dependencia de configuración.
\end{itemize}
No conviene empezar con definiciones. La apertura debe presentar un misterio operacional que no pueda resolverse con una sola señal.

\section{Arquitectura didáctica v9 para 180 minutos}

\subsection*{Bloque 1 --- El incidente misterioso (20 min)}
Mostrar: error rate de checkout elevado, CPU normal y p99 alta. Entregar tres hipótesis posibles y pedir una predicción individual: ``¿qué mirarías después y qué resultado esperás?''.

\textbf{Propósito}: activar conocimientos de clases 1--9 y evidenciar la diferencia entre detectar e interpretar. No resolver todavía.

\subsection*{Bloque 2 --- Logs como hechos contextualizados (30 min)}
Comparar texto libre con un evento estructurado. Construir con el curso los campos mínimos: tiempo, servicio, operación, trace ID, entidad de negocio, resultado, duración e intento. Introducir privacidad, secretos, retención y el riesgo de usar los logs como base de datos accidental.

\textbf{Microdemostración}: dar diez líneas de texto ambiguas y luego el mismo evento en JSON. Pedir que encuentren una operación duplicada.

\subsection*{Bloque 3 --- Métricas y experiencia del usuario (35 min)}
Presentar Golden Signals, RED y USE sobre el mismo checkout. Calcular qué oculta la media frente a p95/p99. Mostrar una explosión de cardinalidad causada por \texttt{user\_id} y pedir una corrección.

\textbf{Pausa cognitiva}: cada pareja redacta una métrica útil y una etiqueta que nunca debería usarse.

\subsection*{Pausa (10 min)}
Recoger una pregunta anónima. Elegir una que permita aclarar diferencia entre métrica de servicio, de recurso y de negocio.

\subsection*{Bloque 4 --- Tracing y causalidad distribuida (35 min)}
Construir una traza de \texttt{POST /checkout}. Explicar contexto, spans, atributos, links y sampling. Después agregar un evento Kafka para mostrar que la causalidad asincrónica no equivale a una pila RPC.

\textbf{Error deliberado}: eliminar la propagación de contexto en Inventory y pedir al grupo que describa qué se vuelve imposible de reconstruir.

\subsection*{Bloque 5 --- Debugging guiado por hipótesis (30 min)}
Volver al incidente inicial. Seguir el ciclo: síntoma, hipótesis, predicción observable, búsqueda de evidencia, refutación, intervención y verificación. Introducir la cadena DB lenta $\rightarrow$ pool saturado $\rightarrow$ retries $\rightarrow$ 500.

\textbf{Caso realista}: pérdida parcial del pipeline de logs. Preguntar cómo se observa la pérdida de observabilidad.

\subsection*{Bloque 6 --- Taller y cierre (20 min)}
Cada grupo toma un flujo de su Trabajo Final y entrega:
\begin{enumerate}
\item un SLI de usuario;
\item cuatro métricas RED/USE;
\item un esquema de log;
\item una traza con cinco spans;
\item una falla de telemetría y su mitigación;
\item una pregunta diagnóstica que podría responder con esa evidencia.
\end{enumerate}
Cerrar con un exit ticket: ``una señal que detecta, una que localiza y una que confirma''.

\section{Caso conductor completo}
Usar un checkout con API, Auth, Inventory, Payments, Orders y un evento asincrónico hacia Shipping. El incidente comienza con contención en una query de Orders, pero el primer error visible aparece en Payments porque recibe el deadline agotado. Este diseño obliga a distinguir causa inicial, amplificadores y síntoma visible.

Progresión recomendada:
\begin{enumerate}
\item solo métricas agregadas: se detecta, pero no se localiza;
\item se agrega tracing: aparece el camino lento;
\item se agregan logs correlacionados: se confirma el lock y el retry;
\item se interrumpe el collector: se discute resiliencia de la propia observabilidad.
\end{enumerate}

\section{Qué explicar, por qué y cómo verificarlo}
\begin{itemize}
\item \textbf{Observabilidad vs monitoreo}: evita confundir dashboards con capacidad de investigación. Verificar pidiendo una pregunta no prevista.
\item \textbf{Logs estructurados}: permiten correlación y consulta; verificar haciendo reconstruir una operación duplicada.
\item \textbf{Percentiles}: representan colas que la media oculta; verificar con un conjunto de latencias pequeño.
\item \textbf{Cardinalidad}: conecta diseño semántico con costo operacional; verificar pidiendo estimar series producidas por etiquetas.
\item \textbf{Tracing}: reconstruye causalidad, no verdad absoluta; verificar rompiendo propagación de contexto.
\item \textbf{Sampling}: introduce sesgo y costo; verificar comparando head y tail sampling para errores raros.
\item \textbf{Debugging causal}: evita culpar al primer componente rojo; verificar exigiendo una hipótesis refutable.
\end{itemize}

\section{Preguntas socráticas}
\begin{itemize}
\item ¿Qué hecho observamos y qué parte es todavía una inferencia?
\item ¿Qué señal distinguiría entre espera por pool y cómputo lento?
\item ¿Qué request real representa este percentil?
\item ¿Dónde se pierde la causalidad si el contexto no cruza el broker?
\item ¿Qué dato no registrarías aunque facilitara el debugging?
\item ¿Cómo sabrías que el pipeline de logs está perdiendo eventos silenciosamente?
\item ¿Qué observación refutaría nuestra explicación del incidente?
\end{itemize}

\section{Errores previsibles y estrategia de corrección}
\begin{itemize}
\item \textbf{``Más logs resuelven todo''}: pedir costo, privacidad y pregunta concreta que cada campo responde.
\item \textbf{``CPU normal significa servicio sano''}: mostrar espera por I/O o pool saturado.
\item \textbf{``Promedio bajo significa buena latencia''}: hacer visible p99 y su relación con timeouts.
\item \textbf{``El primer span rojo es la causa''}: construir una dependencia upstream que agotó el deadline.
\item \textbf{``Trace ID y request ID son lo mismo''}: distinguir unidad de transporte, operación lógica y saga.
\item \textbf{``Kafka ordena toda la traza''}: recordar orden por partición y causalidad por identificadores.
\item \textbf{``La observabilidad siempre está disponible''}: introducir pérdida de logs y buffering local.
\end{itemize}

\section{Criterios de evaluación formativa}
Una respuesta excelente:
\begin{itemize}
\item separa hechos, hipótesis y evidencia;
\item elige señales según una pregunta operacional;
\item incluye dimensiones y unidades;
\item reconoce sampling, cardinalidad, privacidad y costo;
\item reconstruye una cadena causal completa;
\item propone una forma de refutar su propia explicación.
\end{itemize}
Una respuesta suficiente identifica el papel de logs, métricas y trazas y propone correlación básica. Una respuesta insuficiente enumera herramientas o productos sin explicar qué pregunta resuelven.

\section{Retroalimentación durante el encuentro}
Dar feedback en tres capas: primero reconocer la hipótesis correcta, luego señalar la evidencia que falta y finalmente pedir una condición de refutación. Evitar corregir inmediatamente con la solución completa; la clase gana valor cuando el estudiante experimenta el pasaje desde una intuición plausible hacia una explicación comprobable.

\section{Diseño de experiencia para una clase muy bien valorada}
La percepción positiva debe surgir de aprendizaje visible, no de simplificar el contenido. Para ello:
\begin{itemize}
\item anunciar al comienzo la pregunta que podrán responder al final;
\item reutilizar un único caso conductor, evitando ejemplos desconectados;
\item alternar exposición con predicciones breves cada 10--15 minutos;
\item mostrar una herramienta o salida real, no solo diagramas;
\item cerrar cada bloque con ``qué problema resuelve / qué no resuelve'';
\item reservar tiempo para preguntas anónimas y responder al menos una en profundidad;
\item entregar una hoja-resumen con Golden Signals, RED, USE y flujo de debugging;
\item devolver el exit ticket al inicio de la clase siguiente con dos aprendizajes y una confusión común.
\end{itemize}
Estas prácticas elevan claridad, participación y sensación de progreso sin sacrificar rigor.

\section{Vinculación con el Trabajo Final Integrador}
Cada grupo debe producir un ADR de observabilidad para un flujo crítico: SLI, señales, esquema de contexto, política de sampling, alerta y falla de la telemetría. En la revisión docente, preguntar: ``¿qué incidente concreto podrías diagnosticar con esto y cuál seguiría siendo invisible?''.

\section{Bibliografía y fuentes para preparar la clase}
\begin{itemize}
\item Google, \emph{Site Reliability Engineering}: capítulos sobre monitoring y distributed systems monitoring.
\item OpenTelemetry: modelo de trazas, propagación de contexto y semantic conventions.
\item Brendan Gregg: USE Method; Tom Wilkie: RED Method.
\item Postmortems públicos de proveedores de infraestructura para mostrar cadenas causales y fallas de telemetría.
\item Kleppmann, \emph{Designing Data-Intensive Applications}: capítulos sobre sistemas distribuidos y procesamiento de streams como contexto conceptual.
\end{itemize}

\section{Cierre y puente}
La clase siguiente puede comenzar con esta afirmación: observar una cascada de fallas no evita que ocurra. La pregunta pasa entonces a ser cómo contenerla mediante timeouts, retries correctos, circuit breakers, bulkheads, backpressure y degradación elegante.

\end{document}
